% ============================================================
% ПРИЛОЖЕНИЕ А. Листинги ключевых алгоритмов
% ============================================================
\newpage
\section*{ПРИЛОЖЕНИЕ А}
\addcontentsline{toc}{section}{ПРИЛОЖЕНИЕ А. Листинги ключевых алгоритмов}
\begin{center}
  \textbf{Листинги ключевых алгоритмов расчёта}
\end{center}

\subsection*{А.1. Главный модуль расчёта (calculateBeam.ts)}

\begin{lstlisting}[language=Java,caption={Диспетчер расчёта балки}]
import type { BeamConfig, BeamResult } from '@entities/beam';
import { calculateCantilever } from '../lib/cantilever';
import { calculateSimplySupported } from '../lib/simplySupported';
import { calculateOverhang } from '../lib/overhang';

export function calculateBeam(config: BeamConfig): BeamResult {
  let result: BeamResult;
  switch (config.type) {
    case 'cantilever':
      result = calculateCantilever(config);
      break;
    case 'simply-supported':
      result = calculateSimplySupported(config);
      break;
    case 'overhang':
      result = calculateOverhang(config);
      break;
    default:
      throw new Error(
        'Unknown beam type: ' + (config as BeamConfig).type
      );
  }
  return result;
}
\end{lstlisting}

\subsection*{А.2. Расчёт шарнирно-опёртой балки (simplySupported.ts)}

\begin{lstlisting}[language=Java,caption={Вспомогательные функции дискретизации и интегрирования}]
const N_POINTS = 500;

function linspace(start: number, end: number, n: number): number[] {
  const arr: number[] = [];
  const step = (end - start) / (n - 1);
  for (let i = 0; i < n; i++) arr.push(start + i * step);
  return arr;
}

function trapz(y: number[], x: number[]): number[] {
  const result = new Array<number>(y.length).fill(0);
  for (let i = 1; i < y.length; i++) {
    result[i] = result[i - 1]
      + ((y[i - 1] + y[i]) / 2) * (x[i] - x[i - 1]);
  }
  return result;
}
\end{lstlisting}

\begin{lstlisting}[language=Java,caption={Определение реакций опор (фрагмент)}]
export function calculateSimplySupported(
  config: BeamConfig
): BeamResult {
  const { length: L, loads } = config;
  const EI = BASE_FLEXURAL_RIGIDITY;
  const xs = linspace(0, L, N_POINTS);
  let RA = 0, RB = 0;

  for (const load of loads) {
    if (load.type === 'point') {
      const F = load.value, a = load.position;
      const ra = (F * (L - a)) / L;
      RA += ra;
      RB += F - ra;
    } else if (load.type === 'distributed') {
      const { start, end, w1, w2=w1, distribution } = load;
      const len = end - start;
      const totalF = distribution === 'uniform'
        ? w1 * len : ((w1 + w2) / 2) * len;
      const centroid = distribution === 'uniform'
        ? (start + end) / 2
        : start + len*(w1+2*w2) / (3*(w1+w2+1e-12));
      const ra = (totalF * (L - centroid)) / L;
      RA += ra;
      RB += totalF - ra;
    } else if (load.type === 'moment') {
      const M0 = load.value;
      RB += -M0 / L;
      RA += M0 / L;
    }
  }
  // ... (calculation continues)
}
\end{lstlisting}

\begin{lstlisting}[language=Java,caption={Численное интегрирование для деформаций}]
  // Integration for theta(x) and v(x)
  // Boundary conditions: v(0)=0, v(L)=0
  const MoverEI = M_phys.map((m) => m / EI);
  const thetaRaw = trapz(MoverEI, xs);

  // C1 from condition v(L)=0:
  // v(x) = integral(theta)dx = vRaw + C1*x
  const vRaw = trapz(thetaRaw, xs);
  // v(L)=vRaw[L]+C1*L=0 => C1=-vRaw[L]/L
  const C1 = -vRaw[N_POINTS - 1] / L;

  const theta = thetaRaw.map((t) => t + C1);
  const deformedY = vRaw.map((v, i) => v + C1 * xs[i]);

  const maxDeflection = Math.max(
    ...deformedY.map(Math.abs)
  );
\end{lstlisting}


